\section*{Capítulo \ref{ch:g4:air-a-mixture-of-gases} --- O ar, uma mistura de gases}

\begin{solution}{exo:g4:air-a-mixture-of-gases:1}
Uma mistura. Os dois principais gases dela são o nitrogênio e o oxigênio.
\end{solution}

\begin{solution}{exo:g4:air-a-mixture-of-gases:2}
Cerca de quatro quintos do \omterm{def:g4:air-a-mixture-of-gases:air}{ar} são nitrogênio, cerca de um quinto é oxigênio.
\end{solution}

\begin{solution}{exo:g4:air-a-mixture-of-gases:3}
Cerca de 78 litros de nitrogênio e cerca de 21 litros de oxigênio.
\end{solution}

\begin{solution}{exo:g4:air-a-mixture-of-gases:4}
A garrafa não está vazia: ela está cheia de \omterm{def:g4:air-a-mixture-of-gases:air}{ar}, e o \omterm{def:g4:air-a-mixture-of-gases:air}{ar} preso não deixa
a água entrar.
\end{solution}

\begin{solution}{exo:g4:air-a-mixture-of-gases:5}
Há menos oxigênio (16 partes em 100 em vez de 21) e mais dióxido de
carbono (4 partes em 100 em vez de quase nada).
\end{solution}

\begin{solution}{exo:g4:air-a-mixture-of-gases:6}
O oxigênio nos dois casos: uma chama e o nosso corpo usam, os dois, o
oxigênio do \omterm{def:g4:air-a-mixture-of-gases:air}{ar}.
\end{solution}

\begin{solution}{exo:g4:air-a-mixture-of-gases:7}
78 quadrados são nitrogênio; $100 - 78 = 22$ quadrados não são.
\end{solution}

\begin{solution}{exo:g4:air-a-mixture-of-gases:8}
A chama consome o oxigênio do \omterm{def:g4:air-a-mixture-of-gases:air}{ar} que está no pote. Quando sobra pouco
oxigênio demais em volta da chama, ela se apaga.
\end{solution}

\begin{solution}{exo:g4:air-a-mixture-of-gases:9}
Um quinto de 500 litros: $500 \div 5 = 100$ litros de oxigênio.
\end{solution}

\begin{solution}{exo:g4:air-a-mixture-of-gases:10}
As pessoas na sala inspiram oxigênio e expiram dióxido de carbono. Abrir
as janelas traz de volta \omterm{def:g4:air-a-mixture-of-gases:air}{ar} fresco, com mais oxigênio e menos dióxido de
carbono.
\end{solution}

\begin{solution}{exo:g4:air-a-mixture-of-gases:11}
Três vezes mais \omterm{def:g4:air-a-mixture-of-gases:air}{ar} contém três vezes mais oxigênio: cerca de
$3 \times 10 = 30$ segundos.
\end{solution}
